\section{组织与决策：如何在高不确定性里提高胜率}

\subsection{GPA 决策框架的工程化解读}

访谈中的 GPA（Goal, Priority, Alternative）不是管理口号，而是任务分层法。Goal 要快、Priority 要稳、Alternative 要广，三者缺一不可。对研发组织来说，这对应“方向确定、资源聚焦、方案并行”的节奏管理。

\begin{importantbox}{GPA 落地为日常机制}
\begin{itemize}
    \item Goal：周级目标由核心负责人最终拍板，避免无限讨论
    \item Priority：每次只保留少量主线，其他需求进入排队池
    \item Alternative：鼓励并行方案，小步试错后快速淘汰
\end{itemize}
\end{importantbox}

\subsection{行动优先：先拿信息增量，再做抽象结论}

季逸超强调“先做再说”，核心在于信息论意义上的效率：不执行就没有新信息，纯讨论只是在旧信息上循环推导。Agent 产品迭代尤其如此，很多问题只会在真实任务链路中暴露。

\begin{warningbox}{行动优先不等于粗糙上线}
行动优先的前提是有回滚机制、有日志、有灰度边界。没有这些保护，所谓快速试错会演变成不可控事故。
\end{warningbox}

\subsection{“不做什么”是产能时代最关键能力}

AI 工具让产能上升后，最大风险反而是“方向发散”。访谈里反复出现的纪律是：每天明确不做什么，把精力留给最能形成飞轮的那 1--2 条主线。

\begin{knowledgebox}{常见“伪机会”识别清单}
\begin{itemize}
    \item 看起来热闹但与核心能力不耦合
    \item 需要重资产投入但复用率低
    \item 能带来短期曝光却破坏长期产品一致性
\end{itemize}
\end{knowledgebox}

\subsection{本章小结}

在不确定环境中，组织能力的本质是“提高决策吞吐并降低错误成本”。GPA、行动优先和不做清单，三者形成闭环。


